\section{Verification and interpretation}

The strict verification covers the complete experiment-key sets and checks for duplicates. It reconstructs the declared test origins, confirms target dates and prices against the hashed inputs, and recomputes model and persistence errors, ratios, and gates. An independent calculation checks 19,845,948 saved prediction points across the three protocols. Additional reconstruction confirms the training and validation boundaries, sample counts, and target normalization for all 1,980 supervised conditions. These calculations agree with the saved results and the generated aggregate summaries.

Checkpoint replay covers 72 selected cases across six implementations, two price modes, two feature sets, and three horizons, rotating across the five lists. The replay matches the saved inference outputs, configurations, and feature scalers. This is a stratified check, not an independent retraining or replay of every model. The manifests record a fixed Chronos revision, but the separately recorded cache-file hashes describe the inspected cache rather than a byte-level record captured at the original inference run. Neither arithmetic verification nor a zero-shot label proves historical adjustment-factor availability, a representative asset universe, or absence of pretraining-data overlap.

The earlier technical-feature branch compares 990 conditions, or 17,820 supervised rows, with baseline errors from different test dates. The corrected batch aligns every model with its own baseline and removes labels that reach the next evaluation segment. Other corrections and runtime differences also accompany the rerun, so changes from the old archive do not identify the effect of one repair in isolation. Cross-feature and cross-protocol comparisons still have different usable dates and inputs. We therefore use descriptive summaries and make no row-independent confidence-interval or significance claim.
